\documentclass[10pt,a4paper]{article}
\usepackage[margin=22mm]{geometry}
\usepackage[T1]{fontenc}
\usepackage[utf8]{inputenc}
\usepackage{lmodern,microtype,amsmath,amssymb,amsthm,mathtools,mathrsfs,booktabs,longtable,array}
\usepackage[hidelinks]{hyperref}
\newtheorem{theorem}{Theorem}
\newtheorem{proposition}[theorem]{Proposition}
\newtheorem{lemma}[theorem]{Lemma}
\newcommand{\R}{\mathbb R}
\newcommand{\Sym}{\operatorname{Sym}}
\newcommand{\diag}{\operatorname{diag}}
\newcommand{\cone}{\mathcal K}
\allowdisplaybreaks[2]
\setlength{\emergencystretch}{3em}
\setlength{\parskip}{3pt}
\title{A fixed finite C1 testing family\\and an evaluated rational local pencil}
\author{Constructive extraction from the complete Q222 source identities}
\date{23 October 2026}
\begin{document}
\maketitle
\begin{abstract}
The twelve-root C1 identity admits a source-independent finite testing family
on an explicit source box, with no upper bound on $H=h/\sigma$.
A fixed matrix reserve and two algebraic gate variables replace the
point-dependent TF1 radicals. A complementary six-root identity treats
$\lambda\le0$. Keeping both independent state normals and the complete
retained squares gives an evaluated rational source-cone pencil of total
order $669$, with $120$ blocks, maximum block order $12$, and $2694$
auxiliaries. Both directions of its representation of the specified
\emph{local certificate cone} are proved. The construction does not give
an exact lift of the full eight-coordinate $\mathcal Q_{222}$: global
source and independent testing coverage, including the other C1 and C2
boundaries, remains to be extracted.
\end{abstract}

\section{The result and its exact scope}
Write $\mathbf c=(d,m,C,A,D,\alpha,\beta)$ for the seven canonical
coefficients of the reciprocal parent $F_{\mathbf c}$ in
Section~\ref{sec:source-setup}. Let $\cone_{\rm POS}$ and
$\cone_{\rm NEG}$ be the two fixed Gram cones constructed below on
the state charts $W\ge0,Z\ge0$ and $W\ge0,Z\le0$, respectively.
The names refer to state signs, not to the sign of $\lambda$.
Let $\cone_{\rm SIX}$ be the fixed six-root certificate cone of
Section~\ref{sec:six}. Define
\[
 \cone_{669}=(\cone_{\rm POS}\cap\cone_{\rm NEG})+\cone_{\rm SIX}.
\]
These cones are defined by identities in the full polynomial
$F_{\mathbf c}$ and unrestricted positive semidefinite Gram blocks
on the printed testing vectors. They do not impose the nonlinear
stationary source formulas on their visible coordinates.

\begin{theorem}[Explicit local fixed-family admission]\label{thm:main}
Use the native rational stationary source (P1), with $d=1$, and impose
\[
 0<x,\sigma\le2^{-30},\qquad 0<z\le1,\qquad
 0<b,h<1,\quad r=bz<1,\quad E,M,J>0,\quad D,\beta\ge0.
 \tag{D}
\]
Then $\mathbf c\in\cone_{669}$. There is no restriction on the sign of
$\lambda$ and no upper bound on $H=h/\sigma$.
Every $\mathbf c\in\cone_{669}$ satisfies
\[
 F_{\mathbf c}(u,w;W,X,Y,Z)\ge0
 \quad\text{for } |u|\le2^{-24},\ |w|\le2^{-12},
 \quad (W,X,Y,Z)\in\R^4.
 \tag{T}
\]
The statement extends to complex states by real and imaginary splitting.
All testing axes and zero state coordinates are literal specializations.
Both permitted source faces $D=0$ and $\beta=0$ are retained.
Positive multiples of an admitted source are admitted as well.
\end{theorem}

\begin{theorem}[Evaluated local cone representation]\label{thm:pencil}
The supplied rational matrix file defines a homogeneous
affine pencil $L_{669}(\mathbf c,y)$ such that, for every real
$\mathbf c$,
\[
 \boxed{\quad
 \mathbf c\in\cone_{669}\ \Longleftrightarrow\
 \exists\,y\in\R^{2694}:\ L_{669}(\mathbf c,y)\succeq0.
 \quad}
\]
Its total order is $669$, its $120$ diagonal blocks have order at
most $12$, and its upper triangle has $7716$ nonzero rational
coefficient records. There are no unevaluated matrix entries.
\end{theorem}

\begin{center}
\begin{tabular}{lrrrr}
\toprule
Fixed cone & Blocks & Total order & Largest block & Auxiliaries\\
\midrule
$\cone_{\rm POS}$ &52&268&12&1087\\
$\cone_{\rm NEG}$ &49&240&11&934\\
$\cone_{\rm POS}\cap\cone_{\rm NEG}$ &101&508&12&2021\\
$\cone_{\rm SIX}$ &19&161&10&666\\
$\cone_{669}$ &120&669&12&2694\\
\bottomrule
\end{tabular}
\end{center}
The matrix file is \path{C1_UNIFORM_LOCAL_ORDER669.coo.json}.
The last line includes seven source-splitting variables.
Theorems~\ref{thm:main}--\ref{thm:pencil} concern the independent
testing box (T). Feasibility of this pencil does not prove validity
on the other quantum testing regions. In particular this pencil
cannot yet be polarized as an exact full-$\mathcal Q_{222}$ lift.

\paragraph{Proof route.}
Section~\ref{sec:source-setup} states the complete source and parent
identities. Section~\ref{sec:explicit-twelve-roots} extracts the
twelve-root Gram and every retained square.
Section~\ref{sec:sourcebox} supplies numerical source bounds.
Section~\ref{sec:fixed-principal} replaces TF1 by a fixed algebraic
family. Section~\ref{sec:six} handles the complementary circuit
sign. Section~\ref{sec:matrices} proves both local cone inclusions
and specifies the evaluated files. The final section isolates the
global obligation that remains.

\section{Source formulas and the complete parent}\label{sec:source-setup}
The independent source parameters are $x,\sigma,b,z,h$. Source
parameters and testing coordinates are independent throughout.
Set
\[
 r=bz,\quad Q_b=1-b,\quad\theta=1-h,\quad H=h/\sigma,\quad
 e=xH,\quad E=1-e,\quad K=1+x,
\]
\[
 R_s=1+xr,\qquad M=R_s-xz\theta,\qquad
 J=1+xh-\sigma(1-r),
\]
\[
 P=\sigma+K+\sigma z(1+bx+2xh)
              +\sigma bx^2hz^2+x^2h\theta z.
\]
The native domain has $x,\sigma,z>0$, $0<b,r,h<1$,
$E,M,J>0$, and $D,\beta\ge0$. Put
\[
 \begin{split}
 p&=\frac{z(1+xh)E+(1-r)(1+xhz)}{2P},\\
 l&=\frac{(1+\sigma r)\{eR_s(1+\sigma z)+ME\}}{2P},\qquad
 \kappa=\frac{K(1+\sigma r)(1+\sigma z)}{2MP}.
 \end{split}
\]
The canonical source coefficients, with $d=1$, are
\[
 \begin{gathered}
 m=\kappa x^2zH,\qquad C=\frac{Kz(1+\sigma r)E}{2(1-r)P},
 \qquad A=p/r,\qquad \alpha=l/(zQ_b),\\
 D=\frac{K(1+\sigma z)(1+\sigma r)-JP}{2JP},\qquad
 \beta=\frac{K(1+\sigma r)(1+xhz)-\theta P}{2\theta P}.
 \end{gathered}\tag{P1}
\]
In particular $a=Az=p/b$, $\ell=\alpha z=l/Q_b$,
$k=\sqrt x>0$, and $q=\sqrt z>0$. The symbols $A$ and $\alpha$
here are source coefficients; the capital testing gates $A_t,B_t$
used in some displays below are unrelated.

Two exact source identities will provide the bounds. Define
\[
 H_0=1+zQ_b,\qquad B_s=P+xz(J-\sigma),
\]
\[
 \mathcal H=R_sH_0+z\{(1-r)xh(2+xr)
       -e[2+x(1+h+r\theta)]+Ke^2\}.
\]
Direct multiplication of (P1) gives
\[
 2\theta P\beta=hB_s-\sigma H_0,\qquad
 2P(JD+x\theta\beta)=\sigma^2\mathcal H. \tag{BE}
\]

For an arbitrary canonical vector $\mathbf c$ define the full
four-state parent
\begin{align*}
 F_{\mathbf c}={}&4m(W-wX+uY-uwZ)^2
    +4C(1+w^2)(uW-Y)^2\\
 &+4A(1+u^2)(1+w^2)Y^2
    +(4D+2d)(uwW+uX-wY-Z)^2\\
 &+(4\alpha+2d)(1+u^2)(Y-wZ)^2
    +(4\beta+2d)(1+w^2)(X+uZ)^2\\
 &-2d(1+u^2)(1+w^2)Z^2.
 \tag{P4}
\end{align*}
Thus $F_{\mathbf c}$ is linear in $\mathbf c$. Write $F_w=F_{\mathbf c}$
for the source (P1), with $d=1$.

Here are every source coefficient and every auxiliary form needed
in the construction:
\[
 \begin{gathered}
 g=h-Q_b,\quad \nu=(2D-2\beta\theta+h)/\sigma,\\
 \lambda=(8Dh+4\beta\theta^2+2h^2)/\sigma-4bp\sigma z,\qquad
 \lambda_0=\lambda-2H,\qquad K_A=4\kappa H-\lambda,\\
 B_x=8\kappa e,\qquad
 A_y=8(p-l)+4z\{(1+2m)g+2Dh+2p\sigma\},\\
 R_x=-4\{2Db+(1+2m)g-2\beta\theta\},\qquad
 R_0=-4\{2(C+D)(r-1)+2p+r-1\},\\
 \chi=8(pb+lQ_b)-4h+4eh(2D+1)-8Dh\\
 \hspace{14mm}
     +4z\{4Dbh+4pb\sigma+(1+2m)g^2+2\beta\theta^2\},\\
 k_1=2(1+2m)g^2+8Dhb+8pb\sigma+4\beta\theta^2 .
 \end{gathered}\tag{SP3}
\]
Set $t=u^2$, $\xi_w=w^2$ and
\[
 \begin{gathered}
 N_x=X+uhZ,\quad N_y=(Y-bwZ)/q,\quad
 \Gamma=\sigma Z-uwW,\quad R=wZ/q-uqW,\\
 A_\theta=q(t-x)W,\quad P_\beta=uZ-(x/\sigma)wW,\quad
 X_e=X+ewW,\quad V=uw(t-x)W,\\
 J_C=N_y+bR,\qquad J_\ell=N_y-Q_bR,\qquad R_\gamma=R_0-8pt.
 \end{gathered}
\]
Define
\[
 \begin{split}
 E_{yy}&=2z\{(1+2m)t+(1+2D)\xi_w\},\\
 E_{RR}&=z\{8bD+8bp\sigma+4b\sigma(2C+2D+1)(r-1)
                     +k_1t+2b^2(2D+1)\xi_w\},\\
 E_{yR}&=4z\{t[2Dh+2p\sigma+(1+2m)g]+b(2D+1)\xi_w\\
 &\hspace{22mm}+\sigma(2C+1)(r-1)
                     +2D[1+\sigma(r-1)]+2p\sigma\},\\
 \mathcal E_+&=4\beta(1+w^2)N_x^2+2(2D+1)(uX_e)^2
       +4bpzt\Gamma^2\\
 &\hspace{6mm}+4a[(1+t)(1+w^2)-1]J_C^2
       +4\ell tJ_\ell^2+4Cz w^2J_C^2,\\
 \mathscr S&=\lambda\sigma P_\beta^2+
                 (\lambda x/\sigma)(t-x)(w^2-\sigma z)W^2 .
 \end{split}
\]
The complete polynomial identity is
\begin{align*}
 F_w={}&\mathcal E_++[2(1+w^2)+4mw^2]N_x^2
       +[4a+2z(2C+1)]J_C^2+4\ell J_\ell^2\\
 &+E_{yy}N_y^2+E_{RR}R^2+E_{yR}N_yR+(4D/\sigma^2)\Gamma^2\\
 &-4\nu N_x(u\Gamma)+R_\gamma(N_y+bR)(qw\Gamma)\\
 &-4uqw(2m+2D+1)N_xN_y+uqwR_xN_xR\\
 &+qwB_xN_xA_\theta+uA_yN_yA_\theta
        +K_AA_\theta^2+\chi uA_\theta R+\mathscr S .
 \tag{SP6}
\end{align*}
To verify (SP6), substitute the printed linear forms and collect the
ten quadratic state monomials, then the powers of $u,w$. Substitute
(P1) and clear the positive source denominators; each coefficient is
zero. This is an algebraic identity, with both distinct $\Gamma$
couplings displayed explicitly. The accompanying native-parent
verifier performs this generic rational comparison, not numerical
evaluation or a determinant comparison.

The two transformations
\[
 (u,Y,Z)\mapsto(-u,-Y,-Z),\qquad
 (w,X,Z)\mapsto(-w,-X,-Z)
 \tag{SG}
\]
preserve (P4) for every $\mathbf c$. It therefore suffices to construct
the certificate at $u,w\ge0$. An overall state sign permits $W\ge0$.
Both signs of $Z$ will be included.


\section{The complete twelve-root extraction}
\label{sec:explicit-twelve-roots}

This section gives the complete polynomial identity behind the finite
testing family. All lower-case source parameters in this section are in
the coefficient field. The variables $u,w,W,X,Y,Z$ are independent
testing variables. Put $q=\sqrt z$, $x=k^2$, $Q_b=1-b$, and use the
native stationary C1 source of (P1). In particular
\[
 a=p/b,\quad \ell=l/Q_b,\quad e=xh/\sigma,\quad
 \rho=\lambda x/\sigma,\quad L_0=\lambda\sigma,\quad
 q_0=\lambda k,\quad d_G=4D/\sigma^2,\quad S=d_G+\rho.
\]
Here $\lambda>0$. Set $a_0=k/\sigma$, $c=\max(0,1-K_A)$ and
$K_+=K_A+c\ge1$. No inverse of $D$, $\beta$, or $c$ is used.

The polynomial is the complete parent (P4), with all four state
coordinates independent. Its full source definitions are (P1) and (SP3).


\subsection{The twelve principal roots}

Define
\[
\begin{gathered}
 N_x=X+uhZ,\qquad N_y=(Y-bwZ)/q,\qquad R=wZ/q-uqW,\\
 \Gamma=\sigma Z-uwW,\qquad G_* =\sqrt{d_G}\,\Gamma,\\
 A_\theta=q(u^2-k^2)W,\qquad T_\theta=kq(u-k)W,\\
 U_* =\sqrt{q_0}\,q\sqrt u(u-k)W,\qquad
 K_* =\sqrt{q_0k}\,q(u-k)W.
\end{gathered}
\]
For $u,w,W,Z\ge0$ put
\[
\begin{gathered}
 B_1=\sqrt{L_0}\,uZ,\quad B_2=\sqrt\rho\,uwW,\quad
 B_3=\sqrt\lambda\,k^2qW,\quad B_4=\sqrt\lambda\,kwZ/q,\\
 G=\sqrt\lambda\,k\sqrt{uwWZ},\qquad d_i=B_i-G.
\end{gathered}
\]
Thus
\[
 \psi=(N_x,N_y,R,G_*,A_\theta,T_\theta,d_1,d_2,d_3,d_4,U_*,K_*)^{\mathsf T}
       =A_{\rm src}\phi,
\]
where $A_{\rm src}$ has only source entries and
\[
 \phi=(X,Y,Z,W,uZ,wZ,uW,uwW,u^2W,u^{3/2}W,u^{1/2}W,
             \sqrt{uwWZ})^{\mathsf T}.
\]
This proves the claimed separation-rank bound twelve. In particular,
the fourth-root circuit does not add an independent principal testing
function: $B_1B_2B_3B_4=\lambda^2x^2u^2w^2W^2Z^2$ gives the printed
formula for $G$ exactly.

If $W\ge0$ and $Z\le0$, retain the same signed $B_i$ and put $d_i=B_i$.
The twelfth column of $A_{\rm src}$ is then zero. The full matrix is
printed in \texttt{C1\_ROOT\_MATRIX.json} for both signs.

\subsection{Both state-normal conductors and the exact Gram}

Write $J_C=N_y+bR$ and $J_\ell=N_y-Q_bR$. The two linear forms
\begin{align*}
 \mathfrak U={}&\frac{\sqrt{d_G}uG_*
 +\sqrt\rho\{k(d_1-d_2)-(u-k)(d_2-d_3)\}
 -\rho\sqrt\sigma T_\theta}{S},\\
 \mathfrak V_0={}&\frac{R_0\sqrt{d_G}qwG_*+R_0\rho\sigma q^2R
 -R_0\sqrt\rho(qw+q^2\sqrt\sigma)(d_2-d_3)}{S}\\
 &+\frac{R_0q_0q^2+(R_0\sqrt\lambda)\sqrt\rho qw}{S}T_\theta,
 \qquad \mathfrak V=\mathfrak V_0-8pqwu\mathfrak U
\end{align*}
satisfy, exactly,
\[
 \mathfrak U=u\Gamma,\qquad
 \mathfrak V=(R_0-8pu^2)qw\Gamma.
\]
Both complete normal terms are therefore retained. These formulas
are also literal when $D=0$.

Let
\[
 E_{RR}^0=z\{8bD+8bp\sigma+4b\sigma(2C+2D+1)(r-1)\},
\]
and let $E_{yR},\nu,R_0,R_x,B_x,A_y,\chi$ have the explicit source
definitions in (SP3). Put
\[
 \varepsilon_C=2/\lambda+u+k,\qquad
 \varepsilon_G=2a_0u+(u+k)/L_0.
\]
The exact principal Gram is the following quadratic form in the
twelve formal coordinates $\psi$:
\begin{align*}
 \psi^{\mathsf T}\mathcal Q\psi={}&
 2N_x^2+J_C^2+J_\ell^2+G_*^2
 +\frac{K_+}{2}(A_\theta^2+T_\theta^2)
 +\frac12\sum_{i=1}^4d_i^2+U_*^2+\frac12K_*^2\\
 &+\left(E_{RR}^0-\frac{\lambda x}{2}-\frac{cd_Gu^2}{S}\right)R^2
 +E_{yR}N_yR-4\nu N_x\mathfrak U+J_C\mathfrak V\\
 &-4uqw(2m+2D+1)N_xN_y+uqwR_xN_xR
 +qwB_xN_xA_\theta+uA_yN_yA_\theta+\chi uA_\theta R\\
 &+\sqrt{q_0u}\,U_*R
 -\frac cS(\sqrt{a_0u}\,U_*+\sqrt{a_0k}\,K_*)^2
 -\frac{c\varepsilon_G}{S}G_*^2
 -\frac{cd_G\varepsilon_C}{S}\sum_{i=1}^4d_i^2.
\end{align*}
In $s=\sqrt u,w$, this matrix has just the testing monomials
\[
 1,\quad w,\quad w^2,\quad s,\quad s^2,\quad s^2w,
 \quad s^4,\quad s^4w.
\]
The eight evaluated symmetric coefficient matrices, comprising 104
stored upper-triangle entries, are in
\texttt{C1\_PRINCIPAL\_GRAM.json}.

\subsection{The nineteen elementary retained roots}

Define
\[
 X_e=X+ewW,\qquad F_A=u(Z-a_0wW),
\]
\[
 k_1=2(1+2m)(h-Q_b)^2+8Dhb+8pb\sigma+4\beta(1-h)^2.
\]
The following nineteen roots are retained without any mixed square
root. Every displayed radical is purely source-side.
\begin{align*}
 &2\sqrt\beta N_x,\quad2\sqrt\beta wN_x,
 \quad\sqrt{2(2D+1)}uX_e,\quad2\sqrt{bp}\,qu\Gamma,\\
 &2\sqrt a\,uJ_C,\quad2\sqrt a\,wJ_C,
 \quad2\sqrt a\,uwJ_C,\quad2\sqrt\ell\,uJ_\ell,
 \quad2\sqrt C\,qwJ_C,\\
 &\sqrt{2+4m}\,wN_x,\quad
 \sqrt{4a+2z(2C+1)-1}\,J_C,\quad\sqrt{4\ell-1}\,J_\ell,\\
 &\sqrt{2(1+2m)}\,quN_y,\quad
 \sqrt{2(1+2D)}\,qwN_y,\quad
 \sqrt{k_1}\,quR,\quad\sqrt{2(2D+1)}\,qbwR,\\
 &\sqrt{L_0/2}\,F_A,\quad
 \sqrt{K_+/2}\,qu(u-k)W,\quad
 \sqrt{K_+k}\,q\sqrt u(u-k)W.
\end{align*}
The final pair represents $K_+(A_\theta^2-T_\theta^2)/2$.
The native source bounds give $4a>1$ and $4\ell>1$, so their
radicands are nonnegative. Each root is a source-linear combination
of at most three fixed testing monomials.

\subsection{Eight angular roots, and the negative-state correction}

Let
\[
 \mathscr C=\sum_{i=1}^4B_i^2-4\lambda xuwWZ,
 \qquad
 \Delta_A=\varepsilon_C\mathscr C+u^2R^2
                   +\varepsilon_G\Gamma^2-A_\theta^2.
\]
For $W,Z\ge0$ put
\[
 L=\sqrt{q_0\sigma}\,F_A+\frac{k}{\sqrt{q_0\sigma}}\Gamma,
 \quad K_Z=k^2qW-kwZ/q,\quad H_Z=\sqrt{wWZ}.
\]
The elementary angular identity used here is, for $U,V,Z_*\ge0$
and $a_*>0$,
\[
 \Psi=1+U^2V^2+U^2Z_*^2+V^2Z_*^2-4UVZ_*,
 \quad T_*=U^2-1,\quad Q_*=VZ_*-U,
\]
\[
 R_*=Z_*-UV,\qquad G_*=U(V-Z_*),
\]
\begin{align*}
 &[2+a_*(U+1)]\Psi+U^2Q_*^2+
       [2U+(U+1)/a_*]R_*^2-T_*^2\\
 &=(T_*+UQ_*)^2+2Q_*^2
       +(U+1)(\sqrt{a_*}G_*-R_*/\sqrt{a_*})^2\\
 &\qquad+a_*(U+1)(1-VZ_*)^2+
                  2a_*(U+1)VZ_*(U-1)^2 .
 \tag{US}
\end{align*}
Multiplication verifies it. For $W\ne0$, substitute
\[
 U=u/k,\quad V=w/(\sqrt{\sigma}q),\quad
 Z_*=\sqrt{\sigma}Z/(kqW),\quad a_*=q_0
\]
and multiply by $(k^2qW)^2$. The resulting formulas below no longer
divide by $W$ or $Z$, so their polynomial identity includes those
zero coordinates.

Then $\Delta_A$ is the sum of the squares of exactly these eight roots:
\begin{align*}
 &A_\theta+uR,\qquad\sqrt2kR,\qquad
 L,\qquad\sqrt{u/k}\,L,\\
 &\sqrt{q_0}K_Z,\qquad\sqrt{q_0u/k}K_Z,
 \qquad\sqrt{2q_0k}(u-k)H_Z,
 \qquad\sqrt{2q_0u}(u-k)H_Z.
\end{align*}
This follows by multiplying out (US);
the displayed formulas no longer divide by a state coordinate.
Multiply these roots by the pure-source factor $\sqrt{cd_G/S}$.

If $W\ge0,Z\le0$, first use the eight formulas with $|Z|$ in place
of $Z$, and append the three roots
\[
 2k\sqrt{4+2q_0+q_0^{-1}}\sqrt u H_{|Z|},\quad
 2\sqrt k\sqrt{2q_0+2+q_0^{-1}}uH_{|Z|},\quad
 2u\sqrt u H_{|Z|}.
\]
Their squared sum is exactly the increase in $\Delta_A$ when the
positive state coordinate is replaced by its negative. This proves
the eleven-root formula on the negative-state chart.

\subsection{The three circuit roots and their fixed testing span}

On the positive-state chart set
\[
 a_t=u^{1/4},\quad b_t=w^{1/4},\quad
 \omega=W^{1/4},\quad\zeta=Z^{1/4},\quad
 \gamma=\sqrt{\sqrt\lambda\,k}.
\]
The six testing monomials are
\[
\begin{gathered}
 f_1=a_t^3b_t\omega\zeta^3,\quad
 f_2=a_t^3b_t^3\omega^3\zeta,\quad
 f_3=a_tb_t\omega^3\zeta,\quad
 f_4=a_tb_t^3\omega\zeta^3,\\
 f_5=a_t^3b_t^2\omega^2\zeta^2,\qquad
 f_6=a_tb_t^2\omega^2\zeta^2.
\end{gathered}
\]
The retained $\mathscr C_+=\mathscr C-\sum d_i^2$ is the sum of
the squares of
\begin{align*}
 &\sqrt2\gamma(L_0^{1/4}f_1-\rho^{1/4}f_2),\\
 &\sqrt2\gamma\{(\sqrt\lambda k^2q)^{1/2}f_3
                    -(\sqrt\lambda k/q)^{1/2}f_4\},\\
 &2\gamma\{(\lambda k)^{1/4}f_5-(\lambda k^3)^{1/4}f_6\}.
\end{align*}
On the negative-state chart there is instead the single root
$2\sqrt\lambda k\sqrt{uwW|Z|}$.

The common scalar coefficient retained in the parent is
\[
 q_C=\frac12-\frac{cd_G}{S}(2/\lambda+k+u).
\]
On $0\le u\le\delta^2$, with $u=\delta^2A_t^8$ in the scaled
testing parametrization, write $A_C=cd_G\delta^2/S$ and
$q_C^0=\frac12-cd_G(2/\lambda+k)/S$. Exactly
\[
 q_C=(q_C^0-A_C)+A_C(1-A_t^8),\qquad
 1-A_t^8=\xi^2(1+A_t^4),\quad\xi^2=1-A_t^4.
\]
Thus the three fixed multipliers $1,\xi,\xi A_t^2$ suffice for
each circuit root. There is no point-dependent choice of a radical
or a bounding constant.

\subsection{Exact assembly and boundary specialization}

The principal quadratic form, the nineteen elementary squares,
the weighted angular squares, and $q_C\mathscr C_+$ sum to $F_w$.
This identity was independently verified first against the complete
SP6 expression and then against the original P4 parent after the
generic rational P1 substitution. In the latter check all seventeen
nonzero pre-substitution coefficients in the independent polynomial
variables $W,X,Y,Z,u,w$ vanish identically.

All final testing functions are polynomial in quarter-angle and
quarter-state coordinates. Consequently $u=0,w=0,W=0,Z=0$ are
literal specializations, with no limiting SOS argument. The signs
of the quarter-state parameters do not matter: each circuit root
has one fixed parity in each parameter, so its square is invariant.
The testing sign gauges
$(u,Y,Z)\mapsto(-u,-Y,-Z)$ and
$(w,X,Z)\mapsto(-w,-X,-Z)$ restore the other angle signs. An overall
state sign makes $W\ge0$; the two state charts then cover all real
states. Real and imaginary parts restore complex states.

The complete retained coefficient matrices are in
\texttt{C1\_RETAINED\_ROOTS.json}: the positive-state chart has
33 distinct retained testing monomials and 30 roots before the
circuit coefficient is split; the negative-state chart has 27
monomials and 31 roots. Every retained root has separation rank at
most three. These are local certificates for the printed source
and testing domains; none of these identities asserts that this
one chart is an exhaustive global Q222 testing family.

\section{A fully evaluated source box, with arbitrary \texorpdfstring{$H$}{H}}
\label{sec:sourcebox}
The pointwise estimates in the uniform C1 proof can be replaced by numerical
inequalities. Set
\[
 \epsilon=2^{-30},\qquad \delta=2^{-12}.
\]
Use the native source formulas (P1), the definitions (SP3), and impose
\begin{equation}\label{eq:sourcebox}
 0<x,\sigma\le\epsilon,\quad 0<z\le1,\quad 0<b,h<1,\quad
 r=bz<1,\quad E,M,J>0,\quad D,\beta\ge0.
\end{equation}
There is no lower bound on $z,b,1-b,h,1-h$, and no upper bound on
$H=h/\sigma$. The literal faces $D=0$ and $\beta=0$ are admitted.
The opening variables $x,\sigma,z$ remain positive in this chart; a
literal zero opening is not obtained by cancelling a zero factor.

Here and below every displayed inequality holds in the source real
closed field. Put $k=\sqrt x$, $q=\sqrt z$, and $e=xh/\sigma$.
The following estimates will be deliberately conservative.
\begin{align}
 &1\le P\le1+4\epsilon,\quad
  1-\epsilon\le J,M\le1+\epsilon,\quad
 1\le B_s\le1+6\epsilon,\quad 1\le H_0\le2,\label{eq:bound1}\\
 &0<e\le6\epsilon,\quad
 |\mathcal H-H_0|\le32\epsilon,\quad
 0\le D\le3\sigma^2,\quad 0\le\beta(1-h)\le h,\label{eq:bound2}\\
 &\tfrac13\le p,l,\kappa\le2,\quad\kappa\le1,
 \quad 0<m\le6\epsilon^2,\quad |C(r-1)|\le1,\label{eq:bound3}\\
 &|\nu|\le4,\quad |\lambda_0|\le5,\quad
 |K_A-\{2(1-h)H_0-2h(1-r)\}|\le1024\epsilon,\label{eq:bound4}\\
 &|R_0|\le40(x+\sigma),\quad |B_x|\le48\epsilon,
 \quad |A_y|,|\chi|\le64,\quad |R_x|\le24.\label{eq:bound5}
\end{align}

\paragraph{Proof of the bounds.}
For the proof one may first replace $2^{-30}$ by any positive real
$\epsilon\le2^{-12}$. The formula for $P$ is a sum of nonnegative
terms. Its nonconstant part is at most
$x+2\sigma+3x\sigma+x^2\sigma+x^2\le4\epsilon$.
The formulas for $J,M,B_s$ give \eqref{eq:bound1} directly; in particular
$J-\sigma>0$.

Initially $0<e<1$ follows from $E>0$. Dropping the negative term in
$\mathcal H$ gives $\mathcal H\le4$. The source identities
\[
 2(1-h)P\beta=hB_s-\sigma H_0,\qquad
 2P\{JD+x(1-h)\beta\}=\sigma^2\mathcal H
\]
imply
\[
 eB_s\le xH_0+\sigma\mathcal H\le2x+4\sigma.
\]
Thus $e\le6\epsilon$. Expanding $\mathcal H-H_0$ and taking absolute
values gives
\[
 |\mathcal H-H_0|\le5x+3e+2e^2\le32\epsilon.
\]
The second source identity, with both summands nonnegative, gives
$D\le3\sigma^2$. The first gives
$\beta(1-h)\le hB_s/(2P)\le h$. The numerator of $p$ is at least
$H_0-e\ge1-6\epsilon$; the numerator of $l$ is at least
$M(1-e)$. These inequalities, together with \eqref{eq:bound1}, give
the lower bounds $p,l\ge1/3$. The corresponding upper bounds and
$1/3\le\kappa\le1$ follow by bounding each positive numerator factor.
Now $m=\kappa xze$ and the pole-free expression for $C(1-r)$ prove
the rest of \eqref{eq:bound3}.

Use the exact identities
\begin{align*}
 \nu&=2D/\sigma+\{H_0-ez(J-\sigma)\}/P,\\
 \lambda_0&=2e(1-h)z(J-\sigma)/P
       -2(1-h)H_0/P+8Dh/\sigma-4bp\sigma z.
\end{align*}
They imply $|\nu|\le4$, $|\lambda_0|\le5$, and
\[
 |\lambda_0+2(1-h)H_0|\le72\epsilon.
\]
To bound $K_A$ uniformly in $H$, write
\[
 N_\kappa=(1+x)(1+\sigma r)(1+\sigma z)-MP
            =xF_x+\sigma F_\sigma.
\]
Here $F_x=(N_\kappa-N_\kappa|_{x=0})/x$ is a polynomial with sum of
absolute coefficients $37$ after $r=bz$, and
$F_\sigma=r-1+\sigma rz$. All its arguments belong to $[0,1]$.
The exact pole removal is
\[
 H(2\kappa-1)=\frac{eF_x+hF_\sigma}{MP}.
\]
Since $|(MP)^{-1}-1|\le12\epsilon$ and $(MP)^{-1}\le2$,
\[
 |H(2\kappa-1)+h(1-r)|\le458\epsilon.
\]
Combining this with $K_A=2H(2\kappa-1)-\lambda_0$ proves
\eqref{eq:bound4}, since $2\cdot458+72<1024$.

For $R_0$, eliminate the apparent $C$ pole explicitly:
\begin{align*}
 -PR_0/4={}&z(1-e)\{1+xh-(1+x)(1+\sigma r)\}\\
 &+(1-r)\{1+xhz-P\}-2DP(1-r).
\end{align*}
Taking absolute values gives $|R_0|\le40(x+\sigma)$.
The remaining bounds in \eqref{eq:bound5} follow immediately from
the definitions (SP3), $|h+b-1|\le1$, and \eqref{eq:bound2}--\eqref{eq:bound3}.

\paragraph{The positive-$\lambda$ branch.}
Suppose additionally $\lambda>0$ and put
\[
 S=d_G+\rho,\quad d_G=4D/\sigma^2,\quad
 \rho=\lambda x/\sigma,\quad L_0=\lambda\sigma,\quad q_0=\lambda k.
\]
The exact identity
\begin{align*}
 S={}&2\mathcal H/P+d_G(1-J+2hx)+2eH_0/P\\
 &-2e^2z(J-\sigma)/P-4bpxz
\end{align*}
and the preceding bounds give
\begin{equation}\label{eq:positivebounds}
 1\le S\le10,\quad 0\le d_G,\rho\le10,\quad
 0<L_0\le3,\quad 0<q_0\le6,\quad \lambda x\le10\epsilon.
\end{equation}
For example, the negative terms in the expression for $S$ have
magnitude at most $144\epsilon^2+8\epsilon$; its leading term is
at least $2(1-32\epsilon)/(1+4\epsilon)$. This gives $S\ge1$.
The upper bound follows by dropping the two negative terms and
using $d_G\le12$ before improving it to $d_G\le S\le10$.
Finally $q_0^2=L_0\rho\le30$.

The remaining potentially large expression is controlled without
bounding $\lambda$ itself:
\[
 (x+\sigma)^2H=xe+2\sigma e+\sigma h\le2\epsilon,
 \qquad (x+\sigma)^2\le4\epsilon^2.
\]
Since $\lambda\le2H+5$, it follows that
\begin{equation}\label{eq:Rlambda}
 |R_0\sqrt\lambda|\le100\sqrt\epsilon.
\end{equation}

Let $c=\max(0,1-K_A)$ and $a_0=k/\sigma$. The error bound for
$K_A$ gives $0\le c\le4$. If $c>0$, then $h>1/8$: for
$h\le1/8$ the reference value of $K_A$ is at least $3/2$, and its
error is at most $1/4$. On this branch $L_0=2h+\sigma\lambda_0\ge h>1/8$.
Moreover
\[
 a_0^2(hB_s-\sigma H_0)\le\mathcal H,
 \qquad hB_s-\sigma H_0\ge1/16,
\]
so $a_0\le8$. Consequently
\begin{equation}\label{eq:cbounds}
 c/L_0\le32,\quad ca_0\le32,\quad c/\lambda\le32\sigma.
\end{equation}
When $c=0$, all terms containing these factors are exactly zero;
no positive lower bound for $L_0$ or upper bound for $a_0$ is needed.

\paragraph{Bounds on every coefficient matrix.}
Let $s=\sqrt u$, and write the complete principal Gram from the
source export as
\[
 Q(s,w)=Q_0+\sum_{(i,j)\in\mathcal A}C_{ij}s^iw^j,
 \quad
 \mathcal A=\{(0,1),(0,2),(1,0),(2,0),(2,1),(4,0),(4,1)\}.
\]
Its source reference $B$ has $N_x$ entry $2$, radial block
\[
 \begin{pmatrix}2&2b-1\\2b-1&b^2+(1-b)^2\end{pmatrix},
\]
Gamma entry $1$, two angular entries $(K_A+c)/2$, four circuit
entries $1/2$, and final entries $1,1/2$. There are no other
entries. The exact identity
\begin{align*}
 (N_y+bR)^2+(N_y-(1-b)R)^2
 ={}&\frac{N_y^2+R^2}{3}
 +\frac53\left(N_y+\frac{3(2b-1)}5R\right)^2\\
 &+\frac{1+6b(1-b)}{15}R^2
\end{align*}
proves $B\succeq\frac13I_{12}$ without any choice of whitening matrix.
Here $K_A+c\ge1$ by definition.

For a symmetric matrix use the row norm
$\|A\|_\infty=\max_i\sum_j|A_{ij}|$.
Substitution of \eqref{eq:bound1}--\eqref{eq:cbounds} into the
printed matrices yields
\begin{equation}\label{eq:matrixbounds}
 \|Q_0-B\|_\infty\le2048\sqrt\epsilon,\qquad
 \|C_{ij}\|_\infty\le128\quad((i,j)\in\mathcal A).
\end{equation}
These are finite coefficient bounds, not bounds obtained after a test
has been selected. To make the entrywise audit transparent, the
largest possible $s^2$ diagonal correction is
$c(2a_0+1/L_0)/S\le96$; its adjacent normal coupling is at most
$2\sqrt{10}|\nu|<26$, so its row sum is below $128$.
The three-entry normal rows in the $s^2$ and $s^4w$ matrices are at
most $6|\nu|\sqrt{10}<76$ and $12p\sqrt{10}<76$, respectively.
The $s$ matrix is bounded by $\sqrt6/2+32\sqrt{k}$.
The $w$ matrix is bounded using \eqref{eq:Rlambda}; all its entries
are $O(\sqrt\epsilon)$ with the explicit constants above.
The $w^2,s^2w,s^4$ rows are bounded by $5,32,44$, respectively.
Indeed, the largest $w^2$ row is
$2z(1+2D)(1+b)\le4(1+6\epsilon^2)<5$.
The constant correction is bounded by grouping the pole-free
$C(r-1)$ terms before taking absolute values; its normal rows are
at most $144\sqrt\epsilon$, its radial rows at most $612\epsilon$,
and its circuit rows at most $72\sqrt\epsilon+1000\epsilon$.
All these are below the bound in \eqref{eq:matrixbounds}.

Define the evaluated rational number
\[
 R_*:=128(2\delta+2\delta^2+\delta^3+\delta^4+\delta^5)
 =\frac{563087409156097}{9007199254740992}<\frac1{12}.
\]
Because $2048\sqrt\epsilon=1/16<1/12$, the source matrices
\begin{equation}\label{eq:sourcePSDblocks}
 Q_0-R_*I_{12},\qquad
 \frac{\delta^{i+j}}2(128I_{12}\pm C_{ij})
 \quad((i,j)\in\mathcal A)
\end{equation}
are positive semidefinite. In fact the first is positive definite
with lower bound greater than $1/6$. Thus the fixed testing
construction applies throughout \eqref{eq:sourcebox} with
$\lambda>0$, including both weak source-weight faces.

\section{A fixed algebraic compiler for the twelve-root C1 Gram}
\label{sec:fixed-principal}
This section removes the adaptive testing radicals from the principal
$\lambda>0$ C1 Gram on one explicit testing box. The source admission
condition is finite and algebraic. It is a sufficient condition, and no
claim that it covers every stationary source is made.

Let $\delta=2^{-12}$, $s=\sqrt u$, and write the complete principal Gram
as
\[
 Q(c;s,w)=C_{00}(c)+\sum_{(p,q)\in I} C_{pq}(c)s^pw^q,
 \qquad
 I=\{(0,1),(0,2),(1,0),(2,0),(2,1),(4,0),(4,1)\}.
\]
Here $c$ denotes the source, rather than the angular compensation
scalar; the printed source matrices retain both independent state
normals. Their order is twelve. Define the source reference matrix
$B(c)$ by the diagonal entries
\[
 (2,\,2,\,b^2+(1-b)^2,\,1,\,(K_A+c_A)/2,\,(K_A+c_A)/2,
       1/2,1/2,1/2,1/2,1,1/2)
\]
and the single off-diagonal entry $B_{23}=B_{32}=2b-1$ in one-based
indexing, where $c_A=\max(0,1-K_A)$. Its radial block has determinant
one and trace at most three for $0\le b\le1$. Thus
\[
 B(c)\succeq\tfrac13I_{12}.
\]
An elementary exact verification of its radial reserve is
\[
 (N_y+bR)^2+(N_y-(1-b)R)^2
 =\tfrac13(N_y^2+R^2)
 +\tfrac53\left(N_y+\tfrac35(2b-1)R\right)^2
 +\tfrac{1+6b(1-b)}{15}R^2.
\]
No source-dependent real whitening is used.

For a symmetric source matrix $M$, put
\[
 n(M)=\max_i\sum_j|M_{ij}|.
\]
Every operation here is in the source real closed field, and absolute
values and the finite maximum are semialgebraic functions. Set
\[
 E_0=C_{00}-B,\qquad
 A_{pq}=\delta^{p+q}C_{pq},\qquad r_{pq}=n(A_{pq}).
\]
Our explicit source admission condition is
\begin{equation}\label{eq:finite-admission}
 n(E_0)+\sum_{(p,q)\in I}r_{pq}\le\tfrac16.
\end{equation}
It is imposed together with the printed rational stationary-source
formulas, native positive denominators, $D,\beta\ge0$, and $\lambda>0$.
The condition contains no unknown testing value, no choice of a real
standard part, and no unspecified sufficiently small constant.

The elementary inequality
$2|x_ix_j|\le x_i^2+x_j^2$ proves
$-n(M)I\preceq M\preceq n(M)I$ over every ordered real closed field.
Consequently all the following are positive semidefinite source matrices:
\[
 H_0=C_{00}-\sum_I r_{pq}I\succeq\tfrac16I,
 \qquad H_{pq}^{\pm}=\tfrac12(r_{pq}I\pm A_{pq})\succeq0.
\]
On the fixed testing rectangle $0\le s,w\le\delta$, let
$a=s/\delta$, $b_t=w/\delta$. Direct expansion gives
\begin{equation}\label{eq:fixed-matrix-split}
 Q=H_0+\sum_I\bigl[H_{pq}^+(1+a^pb_t^q)
                         +H_{pq}^-(1-a^pb_t^q)\bigr].
\end{equation}
This is a matrix polynomial identity, not a positivity approximation.


\paragraph{A version with one numerical matrix envelope.}
The same dictionary also applies on the more elementary sufficient domain
\[
 n(E_0)\le\tfrac1{12},\qquad n(C_{pq})\le128\quad((p,q)\in I).
\]
Indeed, put $T_\delta=2\delta+2\delta^2+\delta^3+\delta^4+\delta^5$.
At $\delta=2^{-12}$, rational arithmetic gives
$128T_\delta<1/12$. Replace the source matrices above by
\[
 H_0=C_{00}-128T_\delta I,\qquad
 H_{pq}^{\pm}=\tfrac12\delta^{p+q}(128I\pm C_{pq}).
\]
They are positive semidefinite, with $H_0\succeq I/6$, and the
same matrix identity holds. Here even the matrix envelope is a fixed
numerical constant. Section~\ref{sec:sourcebox} establishes these inequalities throughout
the printed source box. They are numerical source estimates, not
a consequence of pointwise boundedness alone.

\paragraph{Only two gate radicals.}
Write $a=A_t^4$, $b_t=B_t^4$, and introduce
\[
 \xi^2=1-A_t^4,\qquad \eta^2=1-B_t^4.
\]
For every nonnegative integer pair $(p,q)$,
\begin{align*}
 1+A_t^{4p}B_t^{4q}&=1+(A_t^{2p}B_t^{2q})^2,\\
 1-A_t^{4p}B_t^{4q}
  &=\sum_{i=0}^{p-1}(\xi A_t^{2i})^2
    +\sum_{j=0}^{q-1}(\eta A_t^{2p}B_t^{2j})^2.
\end{align*}
Empty sums mean zero. The second identity is a finite geometric sum.
For the seven pairs in $I$, the resulting common multiplier list is
\[
\begin{split}
\mathcal M=\{&1,B_t^2,B_t^4,A_t^2,A_t^4,A_t^4B_t^2,A_t^8,A_t^8B_t^2,\\
             &\eta,\eta B_t^2,\eta A_t^4,\eta A_t^8,
               \xi,\xi A_t^2,\xi A_t^4,\xi A_t^6\}.
\end{split}
\]
In particular there is no new radical for each monomial, and no
$\sqrt{M\pm y(t)}$ depending on an adaptively chosen $M$.

\paragraph{A compact algebraic cover of the testing chart.}
Set
\[
 u=\delta^2A_t^8,\quad w=\delta B_t^4,\quad W=\omega^4,\quad Z=\zeta^4,
 \qquad \omega^8+\zeta^8+X^2+Y^2=1.
\]
Together with the two gate equations, these equations define a compact
real algebraic testing set. They cover exactly the normalized real-state
chart $W,Z\ge0$, $0\le\sqrt u,w\le\delta$. All signs of
$A_t,B_t,\omega,\zeta,\xi,\eta$ are allowed. The common state sign takes
the same-sign state chart $WZ\ge0$ to this chart. In particular
\[
 \sqrt{uwWZ}=\delta^{3/2}A_t^4B_t^2\omega^2\zeta^2
\]
is the nonnegative root for every real point of the cover. Since
$\delta=2^{-12}$, every power $\delta^{j/4}$ appearing in the
substitution is rational.

A primitive rescaling of the twelve testing functions is
\[
\begin{split}
 \chi=(&X,Y,\zeta^4,\omega^4,A_t^8\zeta^4,B_t^4\zeta^4,
 A_t^8\omega^4,A_t^8B_t^4\omega^4,A_t^{16}\omega^4,\\
 &A_t^{12}\omega^4,A_t^4\omega^4,A_t^4B_t^2\omega^2\zeta^2).
\end{split}
\]
The circuit coordinates from the preceding section satisfy
$u^{1/4}=\delta^{1/2}A_t^2$ and $w^{1/4}=\delta^{1/4}|B_t|$.
Their squared identities remain valid when $|B_t|$ is replaced by
$B_t$: every monomial in each root has the same parity in $B_t$.
The same parity argument applies to $\omega$ and $\zeta$.

The original twelve-root vector is $g=R(c)\chi$ for the printed
source matrix, after multiplying its columns by
\[
 (1,1,1,1,\delta^2,\delta,\delta^2,\delta^3,
                 \delta^4,\delta^3,\delta,\delta^{3/2}).
\]
On the negative-state chart use $Z=-\zeta^4$, delete the final
coordinate of $\chi$, and take the first eleven columns of the
negative root matrix $A_-$ from \texttt{C1\_ROOT\_MATRIX.json}.
The exact source-to-feature map is
\[
 R_-=A_-[:,1{:}11]\,
 \diag(1,1,-1,1,-\delta^2,-\delta,\delta^2,\delta^3,
                                  \delta^4,\delta^3,\delta).
\]
Here $[:,1{:}11]$ means the first eleven columns in one-based
notation, equivalently \texttt{A\_minus[:, :11]} in Python.
These three negative column signs are necessary because the
feature vector uses $\zeta^4$, whereas the source root formulas
use the signed state coordinate $Z$.

Combining repeated multipliers in \eqref{eq:fixed-matrix-split} yields
\[
 g^{\mathsf T}Qg
   =\sum_{m\in\mathcal M}(m\chi)^{\mathsf T}G_m(c)(m\chi),
 \qquad G_m(c)\succeq0.
\]
Each $G_m$ is a sum of congruences $R(c)^{\mathsf T}H R(c)$ of the
explicit source matrices $H_0,H_{pq}^{\pm}$. Source-field square
completion therefore gives a finite separated SOS in this one fixed
testing family. Both state normals are preserved because the whole
matrix $Q$ is used, rather than a determinant or a selected restriction.
For the normalized physical slack, multiply every testing function by
$\kappa_t$, with
\[
 \kappa_t^2(1+\delta^4A_t^{16})(1+\delta^2B_t^8)=1.
\]
The denominator is strictly positive on the entire testing set.
This is a pure testing restoration, and the sign of $\kappa_t$ has
no effect on the Gram identity.

\paragraph{Evaluated rational matrices.}
The source files print sixteen feature blocks of length twelve, hence
192 feature slots. Exactly 152 polynomial functions are distinct.
Using only the gate relations, the block products span 280 independent
monomials. No state-sphere reduction is needed: every product is
homogeneous of state degree two, so a relation on the normalized
state chart would extend by radial scaling to an open cone and hence
would already be a polynomial relation. There are 1248 symmetric Gram entries. The fully evaluated
coefficient map has 1872 nonzero integer entries and exact row rank
280. Rational elimination produces an affine pencil of sixteen
$12\times12$ blocks in 280 coefficient variables and 968 auxiliary
variables. It has 2854 nonzero affine entries, with denominators only
one or two. Direct multiplication verifies all 280 coefficient equations.

For precision, this evaluated pencil represents the principal Gram
cone exactly: a coefficient vector is in that cone if and only if an
auxiliary vector makes all sixteen blocks positive semidefinite.
It is not a representation of the full quantum set. The retained
circuit, event, and angular roots belong to additional fixed blocks,
and a global extraction also requires source and testing coverage.

\paragraph{The retained circuit scalar uses the same gates.}
The scalar multiplying the nonnegative circuit remainder is
\[
 q_C=q_{C,0}-r_Cu,\qquad
 q_{C,0}=\tfrac12-\frac{c_Ad_G}{S}\left(\frac2\lambda+k\right),
 \quad r_C=\frac{c_Ad_G}{S}\ge0.
\]
The circuit diagonal of the principal Gram and
\eqref{eq:finite-admission} give $q_{C,0}-r_C\delta^2\ge1/3$.
Exactly,
\[
 q_C=(q_{C,0}-r_C\delta^2)
             +r_C\delta^2\bigl[\xi^2+(\xi A_t^2)^2\bigr].
\]
Thus multiplying the three elementary circuit roots requires only the
fixed multipliers $1,\xi,\xi A_t^2$; no adaptive radical is needed here
either.


\section{The evaluated full-parent local source pencils}
The preceding fixed multiplier construction was applied to every
principal term, and all retained roots from the complete C1 identity
were then restored as their own source-positive Gram blocks.
Each elementary or angular retained root has source/testing rank at
most three. The circuit roots are multiplied by the three fixed
functions $1,\xi,\xi A_t^2$ just derived. No guessed monomial extension
is used in this compilation.

For $Z\ge0$, the resulting fixed cone has 52 blocks, total matrix
order 268, maximum block order 12, and 168 distinct testing functions.
Its coefficient map has 281 independent rows and 1368 Gram entries.
Substituting the full parent, linear in the seven canonical source
coefficients $(d,m,C,A,D,\alpha,\beta)$, and eliminating the
coefficient equations gives an evaluated rational pencil with 1087
auxiliary variables. For $Z\le0$, the common circuit-centre term is
absent from the principal root matrix, so its testing span has eleven
functions. The corresponding fixed cone has 49 blocks, total order
240, maximum block order 11, 143 distinct functions, 221 independent
coefficient equations, and 934 auxiliaries.

Block-diagonal concatenation gives one evaluated rational pencil of
101 blocks and total order 508, in the same seven source coefficients
and 2021 auxiliaries. Every one of its 5326 nonzero upper-triangle
affine entries is printed. Direct rational multiplication verifies all
502 full-parent coefficient equations after elimination.

The equivalence proved by these matrices is exact and limited in
scope: a canonical source belongs to both displayed fixed Gram cones
if and only if this pencil admits a positive-semidefinite completion.
Feasible Grams imply nonnegativity of its full slack on both testing
state charts. Conversely, the uniform source estimates together with
the explicit source Gram assignments prove membership for every
admitted native source. The exact gauges
\[
 (u,Y,Z)\mapsto(-u,-Y,-Z),\qquad
 (w,X,Z)\mapsto(-w,-X,-Z)
\]
preserve the full parent and extend the conclusion to
$|u|\le2^{-24}$, $|w|\le2^{-12}$ and every real state. Applying the
real quadratic identity to real and imaginary state parts covers
complex states. The final polynomial identities specialize directly
to testing axes and zero state coordinates. No state coordinate,
$D$, or $\beta$ is inverted in the final family.

These two pencils are full-state local certificates. They do not
represent the full quantum set, and they do not establish that an
arbitrary globally valid C1 or C2 source belongs to the same fixed
cone. The uncovered source and testing regions remain a separate
constructive obligation.


% Include as a chapter or section fragment. All matrices are printed in the JSON.
\section{An explicit fixed family from the six-root identity}\label{sec:six}

This construction represents the complete reciprocal C1 parent on an
explicit source domain and an independent testing rectangle. It proves
membership in a particular fixed Gram cone on that domain. It makes no
assertion that the domain covers all C1 sources or any of the omitted C2
source families.

\paragraph{Source notation.}
Use (P1), (SP3), $Q=Q_b=1-b$, and the source quantities
\[
 L=\lambda/\sigma,\qquad s=\sqrt{1+L^2},\qquad d_G=4D/\sigma^2.
\]
Within this section $s$ is this pure-source normalizer, not the
testing square root used in the preceding section.


\paragraph{The exact full-state Gram.}
For independent testing coordinates $u,w,W,X,Y,Z$ set
\[
\begin{aligned}
 N_x&=X+uhZ,&N_y&=(Y-bwZ)/q,&R&=wZ/q-uqW,\\
 \Gamma&=\sigma Z-uwW,& A_\theta&=q(u^2-x)W,
 &\widehat V&=suw(u^2-x)W.
\end{aligned}
\]
Let $v=(N_x,N_y,R,\Gamma,A_\theta,\widehat V)^T$.
The exact identity used here is
\[
 F_w=E_{\rm ret}+v^T\!\left(G_0+
                  \sum_{\gamma\in\mathcal A}u^{\gamma_1}w^{\gamma_2}E_\gamma
                          \right)v,\qquad
 \mathcal A=\{00,10,01,11,20,02,21,30\}.
 \tag{SR1}
\]
Here
\[
G_0=\begin{pmatrix}
2&0&0&0&0&0\\
0&2&2b-1&0&0&0\\
0&2b-1&b^2+Q^2&0&0&0\\
0&0&0&d_G&0&L/(2s)\\
0&0&0&0&4\kappa H&0\\
0&0&0&L/(2s)&0&1
\end{pmatrix}.
\]
All unspecified entries of the symmetric matrices below are zero.
Indices follow the six roots just printed, starting at 1.
\[
\begin{array}{c|c|l}
\gamma&(i,j)&(E_\gamma)_{ij}\\\hline
00&(3,3)&z[8bD+8bp\sigma+4b\sigma(2C+2D+1)(r-1)]\\
00&(2,3)&2z[\sigma(2C+1)(r-1)+2D(1+\sigma(r-1))+2p\sigma]\\
00&(4,6)&-sx/2\\\hline
10&(1,4)&-2\nu\\
10&(2,5)&A_y/2\\
10&(3,5)&(\chi+\lambda+s^2\sigma x)/2\\\hline
01&(2,4)&qR_0/2\\
01&(3,4)&bqR_0/2\\
01&(1,5)&qB_x/2\\\hline
11&(1,2)&-2q(2m+2D+1)\\
11&(1,3)&qR_x/2\\\hline
20&(2,2)&2z(1+2m)\\
20&(3,3)&zk_1\\
20&(2,3)&2z[2Dh+2p\sigma+(1+2m)g]\\
20&(4,4)&L\\
20&(5,5)&-s^2\sigma\\
20&(4,6)&s/2\\\hline
02&(2,2)&2z(1+2D)\\
02&(3,3)&2zb^2(2D+1)\\
02&(2,3)&2zb(2D+1)\\\hline
21&(2,4)&-4pq\\
21&(3,4)&-4bpq\\\hline
30&(3,5)&-s^2\sigma/2
\end{array}
\tag{SR2}
\]
In particular, the $N_x\Gamma$, $N_y\Gamma$, and $R\Gamma$
entries retain both complete state-normal couplings. No determinant
substitution has replaced the four-state parent.

For completeness, define $X_e=X+ewW$,
$J_C=Y/q-buqW$, and $J_\ell=(Y-wZ)/q+QuqW$. The retained form is
\[
\begin{aligned}
E_{\rm ret}={}&4\beta N_x^2+(4\beta+2+4m)(wN_x)^2
 +2(2D+1)(uX_e)^2+4bpz(u\Gamma)^2\\
&+[4a+2z(2C+1)-1]J_C^2+(4\ell-1)J_\ell^2\\
&+4a(uJ_C)^2+(4a+4Cz)(wJ_C)^2
       +4a(uwJ_C)^2+4\ell(uJ_\ell)^2.
\end{aligned}\tag{SR3}
\]
All ten coefficients are nonnegative when $4a,4\ell\ge1$.

To prove (SR1), start with the full-parent identity SP6. With
$V=uw(u^2-x)W$, direct multiplication gives
\[
\begin{aligned}
 K_AA_\theta^2+\mathscr S
 &=4\kappa H A_\theta^2+Lu^2\Gamma^2+L\Gamma V
                                       +\lambda uA_\theta R,\\
0&=V^2-\sigma u^2A_\theta^2-u\sigma(u^2-x)A_\theta R
                                      +(u^2-x)V\Gamma.
\end{aligned}
\]
Add $s^2=1+L^2$ times the second identity, replace $V$ by
$\widehat V/s$, and collect the eight testing monomials. Clipping
the two radial squares at one gives (SR3). This proves (SR1)--(SR3)
algebraically at either sign of $\lambda$.

\paragraph{Explicit source guards.}
For a real symmetric matrix set
$\|T\|_{\rm row}=\max_i\sum_j|T_{ij}|$. Choose the fixed rational
constants
\[
 \delta=2^{-12}=1/4096,\qquad \eta=1/56.
\]
The source admission conditions, in addition to the native domain, are
\[
 4a\ge1,\quad4\ell\ge1,\quad d_G\ge1,\quad4\kappa H\ge1,
 \tag{SR4}
\]
\[
 \|E_{00}\|_{\rm row}\le1/12,\qquad
 \delta^{|\gamma|}\|E_\gamma\|_{\rm row}\le1/56
                   \quad(\gamma\in\mathcal A\setminus\{00\}).
 \tag{SR5}
\]
These are finite, explicit inequalities in source-only rational and
algebraic functions; no testing point is used to select them.

The radial block has determinant one and trace at most three, so it
is at least $I_2/3$. An explicit pure-source proof is
\[
 (N_y+bR)^2+(N_y-QR)^2
 =\frac{N_y^2+R^2}{3}
 +\frac{5}{3}\left(N_y+\frac{3(2b-1)}{5}R\right)^2
 +\frac{1+6bQ}{15}R^2.
\]
The $(\Gamma,\widehat V)$ block is at least $I_2/2$: its diagonals
are at least one and $|L/s|\le1$. Thus $G_0\succeq I_6/3$.
Consequently
\[
 B=G_0+E_{00}-I_6/8\succeq I_6/8.
\]
For the seven nonconstant monomials put
\[
 y_\gamma=(u/\delta)^{\gamma_1}(w/\delta)^{\gamma_2},\qquad
 T_\gamma=\delta^{|\gamma|}E_\gamma.
\]
On the fixed independent testing box $0\le u,w\le\delta$ one has
$0\le y_\gamma\le1$. The exact matrix identity is
\[
G_0+\sum_\gamma u^{\gamma_1}w^{\gamma_2}E_\gamma
=B+\frac12\sum_{\gamma\ne00}
 \{(\eta I_6+T_\gamma)(1+y_\gamma)
       +(\eta I_6-T_\gamma)(1-y_\gamma)\}.
\tag{SR6}
\]
The scalar equality $7\eta=1/8$ proves the constant coefficient.
The source matrices $\eta I_6\pm T_\gamma$ are positive semidefinite
by (SR5). All factors $1\pm y_\gamma$ are nonnegative in the testing
field. Factor each source matrix over its own real closed field.
This proves a separated tensor SOS with the \emph{fixed} fourteen
testing radicals $\sqrt{1\pm y_\gamma}$. No adaptive constant $M$
or source-dependent testing function occurs.

\paragraph{A full uniform source rectangle for the nonpositive branch.}
The guards hold for every native source satisfying
\[
 0<x,\sigma\le\varepsilon:=2^{-10}=1/1024,\qquad
 0<z\le1,\qquad \lambda\le0.
\tag{SR7}
\]
No inverse of $\beta$ occurs, so its literal zero face is included.
The proof below forces $D>0$ on this nonpositive branch; the $D=0$
source face is therefore assigned to the complementary positive
branch, rather than being cancelled or omitted.

Here are elementary bounds establishing the claim. From the
definition of $p$, $P\ge1$ and $0<E<1$, one has $p\le1+\varepsilon<2$.
Since every term preceding $-4bp\sigma z$ in $\lambda$ is nonnegative,
\[
 H\le2,\quad h\le2\sigma,\quad e\le2x,\quad
 \beta\theta^2\le2\sigma^2,\quad -8\le L\le0,\quad1\le s\le9.
\]
In particular $\theta\ge1/2$ and $\beta\theta\le4\sigma^2$.
The source formulas give
\[
 1\le P\le1+4\varepsilon,\quad
 1-\varepsilon\le M\le1+\varepsilon,\quad
 1-\varepsilon\le J\le1+\varepsilon,\quad
 B_s=P+xz(J-\sigma)\le1+6\varepsilon.
\]
The exact identity $2\theta P\beta=hB_s-\sigma H_0$, where
$H_0=1+zQ\ge1$, implies $H\ge3/4$. Also
\[
 3/8\le\kappa\le1,\quad1/4\le p\le2,\quad
 1/4\le l\le1,\quad C(1-r)\le1,\quad m\le1/8.
\]
For the lower bound on $l$, use
$2Pl\ge M(1-e)\ge(1-\varepsilon)(1-2\varepsilon)$.
Its upper bound follows from
$eR_s(1+\sigma z)+M(1-e)\le(1+\varepsilon)^2$.
These estimates already prove the first, second and fourth guards in
(SR4).

To bound $D$, use the exact BE4 identity
\[
 2P(JD+x\theta\beta)=\sigma^2\mathcal H,
\]
where
\[
\mathcal H=R_sH_0+z\{(1-r)xh(2+xr)
       -e[2+x(1+h+r\theta)]+Ke^2\}.
\]
The preceding inequalities give $1-5\varepsilon\le\mathcal H\le3$,
and hence $D\le3\sigma^2$. Moreover
\[
 d_G\ge
 \frac{2(1-5\varepsilon)/(1+4\varepsilon)-16\varepsilon}
      {1+\varepsilon}\ge1;
\]
the last inequality follows from
$1-31\varepsilon-68\varepsilon^2\ge0$. Thus (SR4) is proved.

The remaining coefficient estimates are deliberately loose:
\[
 |\nu|\le3,\quad |R_0|\le29,\quad |A_y|\le22,
 \quad|\chi|\le24,\quad |R_x|\le6.
\]
They follow by substituting the displayed bounds in the five explicit
source formulas. For example $|g|\le1$ and
$pb+lQ\le2$; all appearances of $C$ in (SR2) outside whole
positive squares occur as $C(1-r)$, whose bound is one.
Direct row sums in (SR2) then give
\[
 \|E_{00}\|_{\rm row}\le64\varepsilon,\qquad
 \|E_\gamma\|_{\rm row}\le64\quad(\gamma\ne00).
\]
For the constant row, the sharper estimates are
$|(E_{00})_{33}|+|(E_{00})_{23}|\le
42\sigma+12D(1+\sigma)<43\varepsilon$ and
$|(E_{00})_{46}|\le9x/2$.
Thus $64\varepsilon=1/16<1/12$ and
$64\delta^{|\gamma|}\le1/64<1/56$, proving (SR5).
These are finite real inequalities, so they hold over arbitrary
real closed source and testing fields by the same algebraic proof.

\paragraph{The fixed testing span.}
Set
\[
 \phi_{10}=(W,X,Y,Z,uZ,wZ,uW,uwW,u^2W,u^3wW)^T.
\]
The exact $6\times10$ source matrix $T$ satisfying $v=T\phi_{10}$ is
\[
T=\begin{pmatrix}
0&1&0&0&h&0&0&0&0&0\\
0&0&1/q&0&0&-b/q&0&0&0&0\\
0&0&0&0&0&1/q&-q&0&0&0\\
0&0&0&\sigma&0&0&0&-1&0&0\\
-qx&0&0&0&0&0&0&0&q&0\\
0&0&0&0&0&0&0&-sx&0&s
\end{pmatrix}.
\]
The first principal block is $T^TBT$; the other fourteen principal
blocks are
\[
 \tfrac12T^T(\eta I_6\pm T_\gamma)T
 \quad(\gamma\ne00)
\]
on testing vectors $\sqrt{1\pm y_\gamma}\,\phi_{10}$.
All these matrices are positive semidefinite by (SR4)--(SR6).

The independent algebraic testing cover adds fourteen variables
$\zeta_{\gamma,\pm}$ with
\[
 \zeta_{\gamma,\pm}^2=1\pm y_\gamma,\qquad
 b_u^2=u(\delta-u),\qquad b_w^2=w(\delta-w),
\]
and the equations
\[
 W^2+X^2+Y^2+Z^2=1,\qquad
 \kappa_t^2(1+u^2)(1+w^2)=1 .
\]
These eighteen equations in twenty-three real variables define a
compact algebraic set covering the entire test rectangle and state
sphere. Multiply every testing vector by $\kappa_t$ to obtain the
normalized physical slack. Each radical multiplies a whole Gram
block, so either sign of every radical is allowed. Axes and zero
state coordinates are literal specializations. The gauges (SG)
and real/imaginary splitting cover negative angles and complex states.

\paragraph{Preserving the elementary root structure: the primary pencil.}
Keep all fifteen
principal blocks on $\phi_{10}$. Absorb the retained $N_x,J_C,J_\ell$
squares into the first principal block; these three rows already lie
in $\operatorname{span}\phi_{10}$. Group the seven remaining retained
rows only as follows:
\[
\begin{array}{c|c|c}
\text{fixed testing vector}&\text{retained roots}&\text{block order}\\\hline
(wX,uwZ)&wN_x&2\\
(uX,wY,uwW)&uX_e,\ wJ_C&3\\
(uZ,uwY,u^2wW)&u\Gamma,\ uwJ_C&3\\
(uY,uwZ,u^2W)&uJ_C,\ uJ_\ell&3
\end{array}
\]
Each block is the sum of the nonnegative source multiples of the
displayed coefficient-vector outer products. Thus every block is
positive semidefinite over the source field. This compilation has
exactly nineteen blocks, fifteen of order ten and four of orders
$(2,3,3,3)$, with total order $161$ and maximum block order ten.
It preserves the testing support of each existing root rather than
adding arbitrary cross terms among all seventeen base monomials.

The evaluated rational source pencil is
\texttt{six\_root\_structured\_source\_pencil.coo.json}.
Its visible coordinates are $(d,m,C,A,D,\alpha,\beta)$; it has
666 free auxiliary coordinates and 2311 nonzero upper-triangular
rational coefficient records. The common feature-product expansion
has 180 coefficient equations and 846 raw upper-triangular Gram
variables. Exact elimination has rank 180, leaving $846-180=666$
free variables. The matching map and every resulting source and
auxiliary coefficient are supplied and exactly checked.

Here both directions of the local cone representation are algebraic.
An affine-pencil witness supplies positive semidefinite Gram blocks;
the checked coefficient map proves their weighted sum equals the
complete parent. Conversely, any collection of Gram blocks giving
that parent is a solution of the 180 printed linear equations. Take
its entries at the printed free positions as the 666 auxiliaries.
The exact pivot equations then reconstruct every other entry, so the
printed affine pencil represents those same Gram blocks. Thus this
pencil is exactly the specified local fixed Gram cone. This is a
two-sided statement about that cone, not a quantum reverse-inclusion
claim.

An admitted rational source for the general guards is
$\sigma=10^{-4}$, $x=10^{-8}$, $b=1/2$, $z=1$, $h=3/20000$.
It has $\lambda>0$ and is therefore not an example from the
$\lambda\le0$ subdomain. The nineteen source Gram matrices at
this point are printed in the structured fixture, with $331$
nonzero rational upper-triangle entries. The source radical
$s$ cancels from their evaluated entries. This fixture is an
additional exact consistency check; the uniform admission theorem
is proved by the preceding inequalities.

The result is a fixed local testing cone with a proved admission
theorem. It is not a primal affine lift of the full eight-coordinate
set. To obtain such a lift one still needs a finite cover of every
valid source and every independent test by compatible fixed testing
families, including the remaining C1 and C2 boundaries.

\section{Exact affine compilation and both local inclusions}
\label{sec:matrices}
For each fixed chart, write the fixed testing blocks as $f_b(t)$.
Expanding
\[
 \sum_b f_b(t)^{\mathsf T}G_b f_b(t)=F_{\mathbf c}(t)
 \tag{CM}
\]
and reducing the explicitly printed gate equations gives a rational
linear system
\[
 {\cal A}g={\cal B}\mathbf c.
\]
Here $g$ lists the upper-triangular entries of all symmetric $G_b$.
An off-diagonal Gram entry contributes twice its coefficient to
(CM); the matrix file stores the actual entry, not this doubled
coefficient. All coefficients are rational because $\delta=2^{-12}$
has rational quarter powers.

The positive chart has $1368$ Gram variables and $281$ independent
equations. The negative chart has $1155$ variables and $221$
independent equations. The structured six-root cone has $846$
variables and $180$ independent equations. The coefficient maps,
ordered monomials, pivots, and every affine entry are supplied.
Exact rational elimination expresses the pivot entries as affine
functions of $\mathbf c$ and the remaining $1087$, $934$, or $666$
free entries, respectively. The files verify these expressions by
substitution into all the original coefficient equations.

This gives both directions without any positivity inference from a
determinant. If the affine pencil is positive semidefinite, its
blocks are feasible Grams in (CM), hence give a fixed-family
certificate on the entire chart. Conversely, any fixed-family
certificate supplies a solution $g$ of (CM). Choose its free entries
as the auxiliary coordinates. The pivot equations reconstruct its
remaining entries exactly, so the affine pencil equals those same
positive semidefinite blocks.

Because the two state charts must both be certified, their pencils
share the visible coordinates and are intersected:
\[
 L_{508}(\mathbf c,y_+,y_-)=
 \diag(L_{\rm POS}(\mathbf c,y_+),L_{\rm NEG}(\mathbf c,y_-)).
\]
The two circuit-sign \emph{constructions} are alternative sufficient
certificates on their common testing box. We include both by a
source-cone sum. With a free split vector $z_c\in\R^7$ put
\[
 L_{669}(\mathbf c,z_c,y_+,y_-,y_6)
   =\diag\bigl(L_{508}(z_c,y_+,y_-),
                       L_{161}(\mathbf c-z_c,y_6)\bigr).
 \tag{LP}
\]
This proves Theorem~\ref{thm:pencil} from the definition of
$\cone_{669}$. The split is not a positivity-constrained source
parameter: it is simply the visible vector of the first certificate
summand.

For Theorem~\ref{thm:main}, use $z_c=\mathbf c$ and the zero witness
in the six-root cone when $\lambda>0$. The explicit twelve-root
source assignment, numerical reserve, and retained squares provide
both state-chart witnesses. If $\lambda\le0$, use $z_c=0$ and the
six-root source assignment. The larger six-root bound
$x,\sigma\le2^{-10}$ contains (D). For every feasible vector in the
sum, add the two full-parent identities. Their common testing
rectangle is (T), and each identity is nonnegative there. Apply
(SG) and homogeneity in the state to finish the proof.

\paragraph{Reading the actual matrices.}
The primary file uses the tuple
\[
 [\,b,i,j,v,a\,]
\]
to mean that rational number $a$ multiplies coordinate $v$ in
entry $(i,j)$ of block $b$. All indices are zero-based; $i\le j$.
Mirror an off-diagonal entry once. Coordinates $0,\ldots,6$ are
$(d,m,C,A,D,\alpha,\beta)$. The following $2694$ coordinates are
listed in \texttt{auxiliary\_order}; their first seven are the
components of $z_c$. There is no constant matrix. Thus the file is
directly an evaluated collection of rational matrices
\[
 L_{669}(\mathbf c,y)=
 \sum_{j=0}^6c_jA_j+\sum_{k=0}^{2693}y_kB_k.
\]
The source formulas and algebraic roots are used to \emph{prove}
feasibility on (D), not to compute the entries of $A_j$ or $B_k$.

\paragraph{Exact checks supplied with the proof.}
The scripts independently check the generic P1-to-P4 parent
identity, both conductor identities, all retained elementary,
angular, and circuit roots at both state signs, and the final
rational coefficient equations. The last check can be run using
only Python's standard library; symbolic regeneration additionally
uses SymPy. No archived checker is imported, and no numerical
optimization or CHSH comparison is a premise of either theorem.
The negative-state source congruence has the explicit signed
column scaling printed in Section~\ref{sec:fixed-principal}.

\section{What is still missing for the full eight-coordinate set}
The local result resolves a concrete part of the constructive
problem: pointwise bounded-residue arguments have been replaced
by one fixed finite algebraic family on an explicit source and
independent testing box. This does not supply the global Bell
slack cone.

A full extraction still needs a finite algebraic testing cover
$\pi:\mathcal T\to\mathcal P$ onto every physical two-qubit testing
strategy, together with fixed testing vectors $f_b$, for which
every valid affine Bell slack satisfies an identity
\[
 \ell_{\mathbf a}(\pi(t))
   =a_0+\sum_{\nu=1}^8a_\nu\Phi_\nu(\pi(t))
   =\sum_b f_b(t)^{\mathsf T}G_b(\mathbf a)f_b(t),
 \qquad G_b(\mathbf a)\succeq0 .
 \tag{GLOBAL}
\]
The assignment $\mathbf a\mapsto G_b(\mathbf a)$ need not be affine;
the coefficient equations in the Gram variables must be affine.
The fixed functions and their algebraic domains must be independent
of $\mathbf a$. A finite disjoint union of testing charts is allowed,
provided every chart is checked and the resulting fixed cone
contains every valid slack.

For the present construction the uncovered obligations are the
other source regimes, including $z>1$ and the other C1/C2 branches;
literal source openings $x=0$, $\sigma=0$, or $z=0$; and independent
tests outside (T). The state-coordinate zeros and both state signs
\emph{within} (D)--(T) have already been handled. A local positive
reserve cannot be extended across an uncovered source or testing
boundary by merely asserting compactness.

If (GLOBAL) is proved with evaluated finite product maps, physical
evaluation gives the forward inclusion of its moment lift. For
the reverse inclusion, the same Gram identity makes every valid
Bell slack nonnegative on every feasible projected moment vector.
A point outside the closed quantum set has a separating valid
affine slack and is therefore excluded. This is the missing
universal adequacy argument; a determinant identity, one stationary
neighbourhood, or agreement on CHSH does not replace it.

\paragraph{Why collecting the old TF1 radicals is not a finite solution.}
For distinct real numbers $M_1,\ldots,M_N$, the square classes
of $y+M_i$ are independent in
$\R(y)^\times/\R(y)^{\times2}$: the valuation at $y=-M_i$ is odd
in precisely the $i$th factor of any nonempty product. Hence
\[
 [\,\R(y)(\sqrt{y+M_1},\ldots,\sqrt{y+M_N}):\R(y)\,]=2^N.
\]
A fixed finite algebraic function field cannot contain all
adaptive choices. This is an obstruction to keeping the original
radicals unchanged, not a nonexistence result for another fixed
SOS representation. The two-gate construction above is such a
different representation on the explicit box.

Pointwise boundedness also gives no uniform real rectangle without
further estimates. For infinitesimal $x$ and any bounded testing
element $y$, the scalar $1+x\otimes y$ has positive real residue.
But every ordinary rectangle $0<x\le\epsilon$, $y\in\R$ contains
points where $1+xy<0$. The numerical source inequalities proved
here are therefore an essential additional step, not a
consequence of a rank bound alone.

No claim about the reverse inclusion of the existing order-$13744$
or order-$63888$ candidates follows from this construction. The
new evaluated matrices represent the local cone stated in
Theorem~\ref{thm:pencil}, and the full-set extraction remains open.

\section*{Primary source and reproducibility}
The source identities are those of
\texttt{latex/Uniform\_C1\_Certificate.tex} and the complete
companion in the public repository
\[
 \text{\url{https://github.com/anubhav70747/q222-finite-sdp}}
\]
at commit
\texttt{0c1ce2e15ed067a6845464e2468559cea5274f2c}.
The present note supplies the evaluated local constants, fixed
families, and rational elimination in addition to those identities.
The source bundle contains the proof's \LaTeX, all fixed features,
coefficient maps, rational pencils, exact verifiers, and input
hashes. Its README gives one-command verification and regeneration
instructions.



\end{document}

